% Generated draft: requires mathematical and source review.
\documentclass[11pt]{article}
\usepackage[margin=25mm]{geometry}
\usepackage{fontspec}
\setmainfont{Latin Modern Roman}
\usepackage{amsmath,amssymb,mathtools,unicode-math}
\setmathfont{Latin Modern Math}
\usepackage{enumitem,needspace,xurl,hyperref}
\setlength{\parindent}{0pt}
\setlength{\parskip}{5pt}
\setlength{\emergencystretch}{3em}
\newcommand{\N}{\mathbb{N}}
\newcommand{\Z}{\mathbb{Z}}
\newcommand{\Q}{\mathbb{Q}}
\newcommand{\R}{\mathbb{R}}
\newcommand{\C}{\mathbb{C}}
\newcommand{\F}{\mathbb{F}}
\newcommand{\A}{\mathbb{A}}
\newcommand{\PP}{\mathbb{P}}
\newcommand{\E}{\mathbb{E}}
\newcommand{\eps}{\varepsilon}
\newcommand{\dd}{\,\mathrm{d}}
\begin{document}
\section*{UnsolvedMath ID: 7500100}
\textbf{Source problem code:} catalogue entry
\subsection{Definitions and mathematical statement}
Let \(\mathsf{RCA}_0\) be the usual subsystem of second-order arithmetic with basic arithmetic axioms, \(\Delta^0_1\)-comprehension, and \(\Sigma^0_1\)-induction. For a finite set \(X\) of second-order arithmetic sentences, let \(T_X=\mathsf{RCA}_0+X\).

Write \(U\preccurlyeq_{\mathrm{int}}T\) when \(U\) is interpretable in \(T\): an admissible interpretation defines the sorts, domains and relations of \(U\) within \(T\), respecting the logical structure, and \(T\) proves the translations of the axioms of \(U\). This is interpretation, not merely inclusion of theorem sets or implication between the extra axioms. The source does not specify the exact interpretation conventions, such as parameters or whether the base arithmetic must be fixed.

The Friedman--Simpson Interpretability Conjecture asks whether, for every pair of finite sets \(X,Y\) of \emph{published mathematical theorems} expressible in second-order arithmetic,
\[
 T_X\preccurlyeq_{\mathrm{int}}T_Y
 \quad\text{or}\quad
 T_Y\preccurlyeq_{\mathrm{int}}T_X.
\]
The direction is explicit: \(T_X\preccurlyeq_{\mathrm{int}}T_Y\) means that \(T_Y\) interprets \(T_X\). Both directions may hold. No effective procedure deciding the comparison is requested.

The restriction to published or actual mathematical theorems is intentionally informal, as the source background emphasizes. It does not define a syntactic or semantic class of admissible sentences, and different formalizations of a published theorem are not specified. Consequently this is a programmatic comparability conjecture, not a claim about arbitrary finite sets of second-order sentences and not a fully formal standalone sentence until that class and interpretation notion are fixed.
\subsection{Short English statement}
Are extensions of \(\mathsf{RCA}_0\) by finite collections of actual published mathematical theorems always comparable by interpretation? The intended class of theorems is deliberately informal and must not be broadened to arbitrary sentences.
\subsection{Sources}
\begin{itemize}
\item Local UnsolvedMath registry, permanent ID 7500100, source code catalogue entry, title ``Friedman--Simpson Interpretability Conjecture''; proposed year 2011; no individual proposer field supplied; updated 2026-08-06.
\item Montalb\'an, \emph{Open Questions in Reverse Mathematics}, item ``Interpretability conjecture'', distinct from Questions 1--31: \url{https://math.berkeley.edu/~antonio/papers/questionsRM.pdf}. Registry access date 2026-07-29; author-preprint PDF text plus manual normalization. The background describes a public author-hosted preprint/HTML with no explicit redistribution license located; short-statement attribution and release terms require review.
\item The source audit dated 2026-07-29 labels the item apparently open while retaining NEEDS\_REVIEW. The supplied summary attributes the programmatic conjecture to Friedman--Simpson and reports no known resolution, without an additional exact citation or an independently verified status assessment.
\end{itemize}
\end{document}
